\chapter{Reading a problem}

\begin{orient}
\textbf{What this chapter is for.} This chapter teaches the thirty seconds before the first line of algebra. Every problem in this book was built by someone: a constant was chosen, an exponent was inserted, an interval was selected, and the choices were made so that one particular structure would appear and one particular cancellation would happen. The chapter gives you five questions to interrogate an expression with, a protocol that orders those questions, and a vocabulary for the kinds of camouflage a problem-setter uses to hide the structure. At the end you should be able to look at an unfamiliar integral, limit, derivative, differential equation, or puzzle and say, before computing anything, what family it belongs to, what is decoration, and what single move would collapse it.
\end{orient}

\section{The first thirty seconds}

A hard problem is almost never hard because of the algebra. The algebra is downstream. A problem is hard because of what you failed to notice while you were still reading it, and everything after that failure is expensive: three pages of correct manipulation that arrive nowhere, a substitution that is legal and useless, an integration by parts that reproduces the original integrand with the sign unchanged.

The corrective is to treat the statement itself as evidence. Consider what it takes to write down
\[
\int_{0}^{1}\frac{x^{4}(1-x)^{4}}{1+x^{2}}\dd x .
\]
Nobody arrives at this by accident. The exponent $4$ is not $3$ or $5$; the numerator is not $x^{4}(1+x)^{4}$; the interval is not $[0,2]$. Somebody divided $x^{4}(1-x)^{4}$ by $x^{2}+1$, saw that the remainder was the constant $4$, and built the problem backwards from that fact. Once you know a division was intended, the problem is a two-minute exercise; the answer is $\tfrac{22}{7}-\pi$, roughly $0.0012645$, which is the classical proof that $22/7$ overestimates $\pi$. Once you do not know it, no amount of trigonometric substitution will save you.

This is the central claim of the book, and it is worth stating flatly.

\begin{keynote}[The design principle]
A constructed problem carries the fingerprints of its construction. An arbitrary-looking exponent, an interval endpoint that matches a constant inside the integrand, a coefficient that makes a discriminant vanish, a starting configuration whose entries sum to something divisible by the number of entries: these are not coincidences, they are the author's working shown through the paper. Reading them is not cleverness. It is literacy.
\end{keynote}

There is a sharp asymmetry of cost here. Reading a problem well costs thirty seconds and, at worst, wastes them. Reading it badly costs the whole attempt, because a wrong classification is self-reinforcing: having decided that an integral needs parts, you will keep doing parts, and each round will look like progress. The expected value of the first thirty seconds is therefore enormous, and the discipline this chapter asks for is simply to spend them deliberately instead of accidentally.

It also changes what practice means. Practising calculation means doing many problems of a type you already recognise. Practising recognition means looking at many problems you do not yet know how to solve, saying out loud what kind of object each one is, and only then checking. The second is uncomfortable and much faster.

Two objections arrive immediately. The first is that this only works on artificial problems. It is true that a randomly generated integral is usually unsolvable in closed form and that a randomly generated puzzle is usually false; but the problems you will actually be asked, in examinations, in competitions, in this book, and in the applied settings where somebody has already done the modelling, are constructed objects. The second objection is that recognition cannot be taught, only accumulated. This is half right. What can be taught is the list of things to look at, and the order to look at them in. Recognition is then a matter of having looked at enough expressions with the right questions in hand.

\section{Five questions to ask before you compute}

Before touching an expression, ask it five things. They are ordered so that the cheap questions come first, and so that a positive answer to an early one usually makes the later ones unnecessary.

\begin{enumerate}
\item \textbf{What is the domain of the object?} Is this an elementary problem (the answer is built from powers, exponentials, logarithms, and trigonometric functions), a special-function problem (the answer needs $\Gamma$, $\Bfun$, $\dig$, $\Li_{2}$, $\erf$, $\Si$, or a $\zeta$ value), or a discrete problem (the answer is a count, a residue, a parity, or a strategy)? Naming the domain tells you which toolbox to open and, just as importantly, which one to keep shut.
\item \textbf{What symmetry does it have?} Parity in $x$; reflection $x\mapsto a+b-x$ across an interval; the involution $x\mapsto 1/x$ on $(0,\infty)$; periodicity; invariance under swapping two parameters; invariance of a discrete process under a relabelling. Symmetry is the cheapest structure to test and the most likely to be the intended route.
\item \textbf{What is engineered about it?} Hunt for the suspiciously clean feature: a bound that is exactly attainable, an exponent that could have been anything, a constant that makes a numerator divisible by a denominator, a coefficient that makes a quadratic a perfect square, sixteen numbers whose exponents sum to a multiple of sixteen.
\item \textbf{What does it reduce to as raw algebra?} Before any calculus, expand, cancel, divide, apply the Pythagorean and logarithm laws, and rationalise. A surprising proportion of frightening expressions are a short expression wearing a costume.
\item \textbf{What would make it trivial?} Name, explicitly, the one substitution or identity that would collapse the problem to something you can finish in a line. Then ask whether the problem was built to admit exactly that move. This question is run backwards from the answer, and it is the most powerful of the five.
\end{enumerate}

The rest of this section makes each question into a technique. These are meta-techniques: they do not compute anything themselves, they decide what will be computed. Every one of them recurs in every part of the book.

\tech{Simplify as raw algebra first}{core}
\cue An expression whose difficulty is visual: stacked fractions, a numerator and a denominator that share obvious factors, trigonometric squares in both, nested exponentials and logarithms, a polynomial quotient with numerator degree at least that of the denominator.
\method Do the algebra you would do if calculus did not exist. Expand and cancel; divide polynomials until the numerator has lower degree; apply $\sin^{2}+\cos^{2}=1$ and $\sec^{2}-\tan^{2}=1$ to numerator and denominator \emph{separately}; collapse $\eu^{\ln u}$, $\ln(ab)$, and $a^{\log_{a}b}$; rationalise a difference of radicals by its conjugate. Only then classify the problem. The design principle behind this technique is that a setter hides a simple object by multiplying it by something that cancels.

\begin{worked}
$\displaystyle\int_{0}^{1}\frac{x^{4}(1-x)^{4}}{1+x^{2}}\dd x$. Expanding, $x^{4}(1-x)^{4}=x^{8}-4x^{7}+6x^{6}-4x^{5}+x^{4}$, and dividing by $x^{2}+1$ gives
\[
\frac{x^{4}(1-x)^{4}}{1+x^{2}}=x^{6}-4x^{5}+5x^{4}-4x^{2}+4-\frac{4}{1+x^{2}} .
\]
Every term is now a power. Integrating on $[0,1]$: $\tfrac17-\tfrac46+1-\tfrac43+4-\pi=\tfrac{22}{7}-\pi\approx 0.0012645$.
\end{worked}

\begin{trap}
Cancelling a factor without checking that it is nonzero on the interval, and cancelling $\cos^{2}x$ from numerator and denominator on an interval containing $\pi/2$. The cancellation is an identity of functions only where both sides are defined.
\end{trap}

\begin{seealsob}Polynomial long division; trigonometric identity pre-processing; conjugate multiplication.\end{seealsob}

Simplification tells you what the object \emph{is}. The next question is what \emph{kind} of object it is, because the same surface can belong to three different worlds and the worlds have disjoint toolkits. The single most common waste of effort in calculus is hunting for an elementary antiderivative that provably does not exist.

\tech{Name the domain before choosing a tool}{core}
\cue Any expression you are about to attack. Specifically: a logarithm divided by a linear or quadratic factor, an integrand with $\eu^{-x^{2}}$, $\sin x/x$, or $x^{s}$ with $s$ non-integer, a sum of reciprocal squares, a factorial inside a root.
\method Decide, in one sentence, whether the answer lives among elementary functions, among the named special functions, or among integers. Log-over-linear and log-over-quadratic on $[0,1]$ or $(0,\infty)$ almost always land on $\zeta(2)$, $\zeta(3)$, Catalan's constant, or a $\Li_{2}$ value; $x^{s-1}\eu^{-x}$ lands on $\Gamma$; $x^{m-1}(1-x)^{n-1}$ lands on $\Bfun$; a count or a probability lands on an integer or a rational. Then commit to that toolkit.

\begin{worked}
$\displaystyle\int_{0}^{1}\frac{\ln x}{1+x}\dd x$. There is no elementary antiderivative, and hunting for one is the failure mode. The domain is the $\zeta$ family, so expand the geometric series and integrate termwise: $\int_{0}^{1}x^{n}\ln x\dd x=-\frac{1}{(n+1)^{2}}$, hence
\[
\int_{0}^{1}\frac{\ln x}{1+x}\dd x=\sum_{n\geq 0}(-1)^{n}\cdot\frac{-1}{(n+1)^{2}}=-\sum_{m\geq 1}\frac{(-1)^{m-1}}{m^{2}}=-\frac{\pi^{2}}{12}\approx -0.822467 .
\]
\end{worked}

\begin{trap}
Declaring the special-function domain too early. Many integrals with $\ln$ and $\arctan$ in them are elementary after one reflection; check the symmetry question before you reach for a polylogarithm.
\end{trap}

\begin{seealsob}Termwise integration of a Maclaurin series; the log-power moment; Gamma-function moment.\end{seealsob}

Having named the domain, test the cheapest structure. Symmetry costs one substitution to check and, when it is present, usually finishes the problem outright. The specific form to look for is an \emph{involution}: a map of the domain to itself which, applied twice, is the identity. Reflection $x\mapsto a+b-x$ on $[a,b]$, negation $x\mapsto -x$ on $[-a,a]$, and inversion $x\mapsto 1/x$ on $(0,\infty)$ are the three that account for most of the definite integrals in this book.

\tech[Look for the involution]{Look for the involution}{core}
\cue A definite integral whose interval is fixed by an obvious map: $[-a,a]$ under $x\mapsto -x$, $[a,b]$ under $x\mapsto a+b-x$, $(0,\infty)$ under $x\mapsto 1/x$, $[0,\pi/2]$ under $x\mapsto \pi/2-x$. In discrete problems: a process invariant under reversing the order of moves, or a bijection pairing each configuration with a partner.
\method Apply the involution, call the result $I$ as well, and \emph{add}. If $g(x)+g(\sigma(x))$ is constant, or if the two copies differ by a sign, the integral evaluates without an antiderivative. The three outcomes are: the sum is constant (answer is length times constant over two), the sum is zero (answer is zero), or the two are equal (you have gained an identity to feed to a later step).

\begin{worked}
$\displaystyle J=\int_{0}^{\infty}\frac{\ln x}{1+x^{2}}\dd x$. Under $x\mapsto 1/x$ the interval is fixed, $\dd x\mapsto \dd x/x^{2}$, and $\frac{\ln x}{1+x^{2}}\dd x\mapsto -\frac{\ln x}{1+x^{2}}\dd x$. So $J=-J$ and $J=0$. No antiderivative is needed, and none is pleasant.
\end{worked}

\begin{trap}
Forgetting that the involution must fix the \emph{interval}, not merely the integrand. On $(0,1)$ the map $x\mapsto 1/x$ sends the interval to $(1,\infty)$; that is a useful move, but it is a change of region, not a symmetry, and the two pieces must then be recombined.
\end{trap}

\begin{seealsob}King's rule; parity splitting on a symmetric interval; complementary sine and cosine.\end{seealsob}

The involution question is really a special case of a broader habit: asking which features of the problem are load-bearing and which are decoration. A setter who wants to hide a reflection will bury it under an exponent chosen to look forbidding. The counter-move is to replace the forbidding thing by a symbol and watch it fail to appear in the answer.

\tech{Delete the ugly constant}{core}
\cue A constant that has no business being there: an exponent of $\sqrt{2}$ or $\pi$ or $2026$, a coefficient like $1/7$ attached to one term, an evaluation point such as $x=67$, a modulus of $2027$, a nested function like $\cos(x^{99})$ that is doing nothing near the point of interest.
\method Replace the offending constant by a letter, or delete the offending factor entirely, and solve the resulting skeleton problem. If the skeleton has a clean answer and the letter never appears in it, the constant was decoration and the skeleton answer is the answer. If the letter does appear, you have learned that the constant is structural, and you should then ask what identity it completes.

\begin{worked}
$\displaystyle\lim_{x\to 0}\frac{\eu^{x^{3}}-1-x^{3}}{x^{6}\cos(x^{99})}$. Delete $\cos(x^{99})$: it tends to $1$ and is continuous, so it cannot affect the value. The skeleton is $\lim_{u\to 0}\frac{\eu^{u}-1-u}{u^{2}}$ with $u=x^{3}$, which is $\tfrac12$ by the second-order expansion of $\eu^{u}$. The limit is $\tfrac12$; numerically the quotient at $x=0.03$ is $0.5000045$.
\end{worked}

\begin{trap}
Deleting a constant that is structural. The exponent $2026$ paired with a modulus of $2027$ is not noise: it completes a cyclotomic identity, since $\prod_{k=1}^{2026}\bigl(x^{2}-2x\cos\tfrac{k\pi}{2027}+1\bigr)=\frac{x^{4054}-1}{x^{2}-1}$. Test the constant, do not assume it.
\end{trap}

\begin{seealsob}Parameter differentiation; the decoy evaluation point; exponent-independent reflection integrals.\end{seealsob}

The last of the five questions inverts the direction of work. Instead of pushing the given expression forwards, you postulate the shape of the answer and differentiate, expand, or count backwards to see what would produce it. This is how antiderivatives that look impossible are found in one line, and it is how a setter's intentions become visible.

\tech{Work backwards from a clean answer}{core}
\cue An integrand of the form $\eu^{x}\cdot(\text{rational})$, $\eu^{x}(f+f')$, or $\frac{u'v-uv'}{v^{2}}$; a sum whose terms are differences of a shifted quantity; an answer that the problem has told you is an integer, or a nice multiple of $\pi$, or independent of a parameter.
\method Guess the family the answer belongs to, write the general member with undetermined pieces, and differentiate (or difference, or count) it forwards. Matching one or two coefficients then pins it down. The reverse product rule $\int \eu^{x}(f(x)+f'(x))\dd x=\eu^{x}f(x)+C$ and the reverse quotient rule are the two workhorses; for series, the corresponding move is to guess $a_{k}=b_{k}-b_{k+1}$ and solve for $b$.

\begin{worked}
$\displaystyle\int_{0}^{1}\frac{x\,\eu^{x}}{(1+x)^{2}}\dd x$. Guess that the antiderivative is $\eu^{x}$ over something linear, and try $\frac{\eu^{x}}{1+x}$. Differentiating,
\[
\dvop{x}\frac{\eu^{x}}{1+x}=\frac{\eu^{x}(1+x)-\eu^{x}}{(1+x)^{2}}=\frac{x\,\eu^{x}}{(1+x)^{2}},
\]
which is the integrand exactly. So the value is $\frac{\eu}{2}-1\approx 0.359141$.
\end{worked}

\begin{trap}
Guessing a family that is too narrow and concluding that no antiderivative exists. If $\eu^{x}\cdot\frac{P}{Q}$ fails with $P$ constant, allow $P$ to be a general polynomial of the degree of $Q$ before abandoning the guess.
\end{trap}

\begin{seealsob}Reverse product and quotient rule; telescoping sums; engineering the antiderivative $v$.\end{seealsob}

\section{The order of attack}

The five questions become a protocol when you fix their order. The order is chosen so that each step is cheaper than the one after it, and so that a success at any step sends you back to the beginning with a smaller problem.

\begin{enumerate}
\item \textbf{Read the question, not the label.} Decide what is actually being asked for: a value, a closed form, an existence claim, a bound, a strategy, a count. A problem filed under ``derivatives'' whose content is a Vieta relation is a Vieta problem.
\item \textbf{Simplify as raw algebra.} If anything cancels, cancel it, and restart at step 1 with the reduced object.
\item \textbf{Name the domain.} Elementary, special-function, or discrete. Commit.
\item \textbf{Test the symmetries.} Parity, reflection, inversion, periodicity, parameter exchange, and, in discrete problems, invariance of a candidate quantity under the move.
\item \textbf{Identify the decoration.} Replace every arbitrary-looking constant by a letter. Solve the skeleton.
\item \textbf{Ask what would make it trivial.} Name the one substitution or identity, and check whether the problem was built to admit it.
\item \textbf{Only now execute.} Choose the standard technique for the family you have identified, and run it.
\item \textbf{Verify.} By a method independent of the one that produced the answer. This is Chapter 2, and it is not optional.
\end{enumerate}

The first six steps should take under a minute on a problem of ordinary difficulty. If they take longer, that is usually a sign that step 2 was done carelessly.

Two refinements. The protocol is a loop, not a list: every success returns you to step 1 with a new object, and the new object deserves the same interrogation as the old one. A substitution that turns an integral over $(0,\infty)$ into an integral over $[0,\pi/2]$ has handed you a fresh interval with fresh symmetries, and failing to re-ask the symmetry question at that moment is the commonest way to solve half a problem and stall.

The second refinement concerns what to do when nothing triggers. If steps 2 through 6 all come back empty, the problem is probably in the special-function or parameter-method families, where the structure is not visible in the statement but is created by you: insert a parameter where none was given and differentiate in it, or introduce an integral representation of one factor and exchange the order. That is a deliberate escalation, and it belongs after the cheap questions rather than instead of them. It is also the point at which a genuinely hard problem announces itself: hard problems are those where you must manufacture the structure, not merely notice it.

\begin{center}
\begin{tikzpicture}[node distance=6mm and 9mm]
\node[dtest] (q1) {Does anything cancel\\ as raw algebra?};
\node[dleaf, right=9mm of q1] (l1) {Cancel it, then\\ restart at the top};
\node[dtest, below=6mm of q1] (q2) {Is the object fixed by\\ an involution?};
\node[dleaf, right=9mm of q2] (l2) {Reflect and add;\\ split by parity};
\node[dtest, below=6mm of q2] (q3) {Is a constant or exponent\\ doing no work?};
\node[dleaf, right=9mm of q3] (l3) {Delete it and solve\\ the skeleton};
\node[dnode, below=6mm of q3] (q4) {Name the domain, then\\ run its standard ladder};
\draw[dedge] (q1) -- node[dlab,above]{yes} (l1);
\draw[dedge] (q1) -- node[dlab,left]{no} (q2);
\draw[dedge] (q2) -- node[dlab,above]{yes} (l2);
\draw[dedge] (q2) -- node[dlab,left]{no} (q3);
\draw[dedge] (q3) -- node[dlab,above]{yes} (l3);
\draw[dedge] (q3) -- node[dlab,left]{no} (q4);
\end{tikzpicture}
\end{center}

\section{Difficulty is layers, not depth}

It is tempting to rank problems by the sophistication of the technique they need. That ranking is close to useless. A problem requiring Feynman's trick and a Beta-function reflection can be routine, because both steps are visible from the statement. A problem requiring nothing beyond a polynomial division can be brutal, because the division is hidden under four layers of disguise. Difficulty, in a constructed problem, measures the number of layers of camouflage between the surface and the structure.

\begin{keynote}[What difficulty measures]
Depth is how advanced the decisive technique is. Camouflage is how many transformations separate the statement from the form in which that technique is visible. A problem's difficulty tracks camouflage far more closely than it tracks depth. A three-layer problem solved by long division is harder than a one-layer problem solved by contour integration.
\end{keynote}

Camouflage comes in a small number of species, and they are worth naming, because a named thing is a thing you can look for.

\textbf{Exponents that do not matter.} An exponent is inserted precisely because it looks like it must be handled. In $\int_{0}^{\pi/2}\frac{\sqrt{\sin x}}{\sqrt{\sin x}+\sqrt{\cos x}}\dd x$ the square roots are inert: under $x\mapsto \pi/2-x$ the integrand becomes its own complement, the two copies sum to $1$, and the value is $\pi/4$ for the square root, the cube root, or any exponent at all. The rule of thumb: an exponent that appears identically in every occurrence of a function is usually decoration.

\textbf{Parameters that cancel.} A parameter is inserted so that the reader will try to track it. In $\int_{-1}^{1}\frac{x^{2}}{1+\eu^{x}}\dd x$ the exponential looks like the difficulty. Writing $g(x)=\frac{1}{1+\eu^{x}}$, the identity $g(x)+g(-x)=1$ means that reflecting and adding leaves $\int_{-1}^{1}x^{2}\dd x$ over two copies, that is $\int_{0}^{1}x^{2}\dd x=\tfrac13$. The exponential never appears in the answer. Any $g$ with $g(x)+g(-x)=1$ would do.

\textbf{Integrands that are secretly a derivative.} The reverse product and quotient rules generate integrands of arbitrary ugliness from clean antiderivatives. If a rational or exponential integrand has a squared denominator and a numerator of one degree less than you expect, differentiate the obvious candidate before doing anything else. The demonstration in the ``work backwards'' block above is exactly this species.

\textbf{Sums that are secretly telescoping.} $\sum_{k\geq 1}\arctan\frac{1}{k^{2}+k+1}$ looks like it needs a transform. It does not: $\arctan(k+1)-\arctan k=\arctan\frac{1}{1+k(k+1)}$, so the sum telescopes to $\frac{\pi}{2}-\frac{\pi}{4}=\frac{\pi}{4}$. The signature is a term whose denominator is a shifted product, or whose form matches a subtraction identity for $\arctan$, $\ln$, or a partial fraction.

\textbf{Expressions that are secretly constant.} The extreme case of camouflage. A construction with nested trigonometry, a geometric series, and a conjugate logarithm can collapse to $x$: for instance $x\cdot\frac{1+\cos 2x}{2}\cdot\sum_{n\geq 0}\bigl(\frac{\tan^{2}x}{1+\tan^{2}x}\bigr)^{n}=x\cos^{2}x\cdot\sec^{2}x=x$, so its derivative is $1$ everywhere and the elaborate evaluation point in the problem statement is decoration. Before differentiating a monster, evaluate it at two convenient points and see whether it moves.

\textbf{Endpoints chosen to make a symmetry exist.} The most delicate species, because the camouflage is in what is \emph{not} written. An integral over $[0,\pi/4]$ rather than $[0,\pi/2]$, or over $[1,\eu]$ rather than $[1,2]$, is announcing which reflection or substitution was intended: $\theta\mapsto \pi/4-\theta$ in the first case, $x\mapsto \eu/x$ or $u=\ln x$ in the second. When an endpoint is an odd-looking constant, ask what map fixes the interval it defines, and what substitution sends that interval to something standard.

Two closing remarks on camouflage. First, the layers compose: a real problem will hide a telescoping sum inside a substitution inside an algebraic simplification, and peeling one layer is supposed to reveal the next, not the answer. If a first move produces something that is merely a different mess, that is normal; ask the five questions again of the new object. Second, camouflage is detectable in the aggregate even when it is invisible in a single term. Symptoms are: an answer promised to be an integer, a problem that specifies far more constants than the answer could plausibly depend on, and an instruction to evaluate at a point that no honest calculation would prefer.

\section{The protocol in action}

Three problems, from three different parts of the book. Each is solved in a few lines, and in each case the lines are short only because the reading was done first. Every value below has been checked numerically.

\subsection{An integral}

\[
I=\int_{0}^{1}\frac{\ln(1+x)}{1+x^{2}}\dd x .
\]

\emph{Domain.} The logarithm over a quadratic on $[0,1]$ suggests the $\Li_{2}$ world, which would be a long road. Hold that thought.

\emph{Symmetry.} The interval $[0,1]$ is not fixed by any obvious involution, and the integrand has no parity. But $1+x^{2}$ in a denominator is the signature of $x=\tan\theta$, and that substitution maps $[0,1]$ to $[0,\pi/4]$, an interval whose reflection $\theta\mapsto \pi/4-\theta$ \emph{is} an involution.

\emph{Engineering.} The $1$ in $\ln(1+x)$ and the upper limit $1$ are the same $1$. That is the fingerprint: change either and the collapse dies.

Substituting $x=\tan\theta$, with $\dd x=\sec^{2}\theta\dd\theta$ and $1+x^{2}=\sec^{2}\theta$,
\[
I=\int_{0}^{\pi/4}\ln(1+\tan\theta)\dd\theta .
\]
Now reflect. Since $\tan\bigl(\tfrac{\pi}{4}-\theta\bigr)=\frac{1-\tan\theta}{1+\tan\theta}$,
\[
1+\tan\Bigl(\frac{\pi}{4}-\theta\Bigr)=\frac{2}{1+\tan\theta},
\]
so the reflected integrand plus the original is $\ln 2$, a constant. Adding the two copies, $2I=\frac{\pi}{4}\ln 2$ and
\[
I=\frac{\pi\ln 2}{8}\approx 0.2721983 .
\]
The special-function domain was a decoy: the answer is elementary because the reflection was available.

\subsection{A limit}

\[
L=\lim_{n\to\infty}\frac{1}{n}\left(\frac{(2n)!}{n!}\right)^{1/n} .
\]

\emph{Domain.} Factorials with an $n$th root read as ``Stirling''. Stirling would work, and it would be four lines of asymptotic bookkeeping.

\emph{Reduction as raw algebra.} $\frac{(2n)!}{n!}=\prod_{k=1}^{n}(n+k)$, so the whole expression is $\Bigl(\prod_{k=1}^{n}\frac{n+k}{n}\Bigr)^{1/n}$: the $n$th root of a product of $n$ factors, each of which depends on $k$ only through $k/n$. That is a geometric mean, and the $1/n$ prefactor was placed exactly to make it one.

\emph{What would make it trivial.} A logarithm, which converts the geometric mean into an arithmetic mean, which is a Riemann sum. So
\[
\ln L=\lim_{n\to\infty}\frac{1}{n}\sum_{k=1}^{n}\ln\Bigl(1+\frac{k}{n}\Bigr)=\int_{0}^{1}\ln(1+x)\dd x=2\ln 2-1,
\]
and therefore $L=\eu^{2\ln 2-1}=\dfrac{4}{\eu}\approx 1.4715178$. At $n=5000$ the expression equals $1.47162$, converging from above as expected for a concave integrand.

\subsection{A puzzle}

The sixteen numbers $2^{0},2^{1},\dots,2^{15}$ are written on a board. Repeatedly, two of them are erased and replaced by their arithmetic mean and their harmonic mean. The sixteen values converge to a common limit $L$. Find $L$.

\emph{Read the question.} It asks for a terminal state under a repeated move, and the answer is claimed to be independent of which pairs are chosen and in what order. Order-independence is the loudest possible signal that an invariant governs the process.

\emph{Hunt the invariant.} Test the candidates against the move $\{a,b\}\mapsto\{\frac{a+b}{2},\frac{2ab}{a+b}\}$. The count is trivially preserved. The sum is not: it changes unless $a=b$. The product is preserved exactly, because
\[
\frac{a+b}{2}\cdot\frac{2ab}{a+b}=ab .
\]
So $\prod_{i}a_{i}$ is invariant across the whole board.

\emph{Engineering.} The initial product is $2^{0+1+\cdots+15}=2^{120}$, and $120$ is divisible by $16$. The setter chose sixteen consecutive powers precisely so that $120/16$ would be a half-integer with a clean radical.

\emph{Existence of the limit.} The invariant alone constrains a limit that might not exist, so this step is obligatory: since $\mathrm{AM}\geq\mathrm{GM}\geq\mathrm{HM}$, each move weakly contracts the spread between the largest and smallest entries and never moves an entry outside the current range, so the values converge to a common $L$. Then $L^{16}=2^{120}$, giving
\[
L=2^{15/2}=128\sqrt{2}\approx 181.0193 .
\]
Simulating four hundred thousand random moves from the given start indeed drives every entry to $181.01934$.

\section{Recognition matrix}
\begin{recog}{@{}p{0.31\textwidth}p{0.35\textwidth}p{0.26\textwidth}@{}}
\rowcolor{mist}\mh{If you see} & \mh{First move} & \mh{If that fails} \\
\midrule
\endhead
A quotient with $\deg(\text{num})\geq\deg(\text{den})$ & Divide; read the remainder & Partial fractions \\
Integral over $[-a,a]$ & Split into even and odd parts & Reflect and add \\
$f(x)+f(a+b-x)$ in a denominator & Reflection $x\mapsto a+b-x$ & Parameter differentiation \\
Integral over $(0,\infty)$ with $x$ and $1/x$ balanced & Inversion $x\mapsto 1/x$ & Glasser or Cauchy-Schl\"omilch \\
An exponent that is $\sqrt2$, $\pi$, or a big integer & Replace it by $p$; expect it to cancel & Test $p=1$ numerically \\
$\ln$ over a linear or quadratic factor on $[0,1]$ & Expand geometrically; expect $\zeta$ or $\Li_{2}$ & Feynman in a parameter \\
Factorials or $n!$ under an $n$th root & Take logs; look for a Riemann sum & Stirling \\
$\frac{1}{n}\sum_{k=1}^{n}f(k/n)$ in any disguise & Recognise the Riemann sum & Squeeze between integrals \\
A product of many factors, derivative wanted & Logarithmic differentiation & Check whether the product telescopes \\
An evaluation point of $0$, $1$, or $\pi/2$ & Find the factor that vanishes there & Expand to first surviving order \\
Numerator one degree below a squared denominator & Reverse quotient rule & Parts with $\dd v=\dd x$ \\
$\eu^{x}\bigl(f(x)+f'(x)\bigr)$ & Reverse product rule: answer is $\eu^{x}f(x)$ & Tabular integration \\
Term of the form $b_{k}-b_{k+1}$ & Telescope & Abel summation \\
$\dv{y}{x}=F(y/x)$ & Substitute $v=y/x$ & Test exactness \\
$\dv{y}{x}+P(x)y=Q(x)$ & Integrating factor $\eu^{\int P}$ & Variation of parameters \\
A move replacing two items by two items & Test sum, product, parity, sum of squares & Monovariant plus termination \\
``Prove this is impossible'' & Colour, or take residues mod $k$ & Extremal principle \\
A statement about its own truth & Encode as a biconditional & Truth table on $2^{n}$ rows \\
An answer claimed independent of a parameter & Set the parameter to its easiest value & Differentiate in the parameter \\
A limit in two variables & Test the paths $y=mx$ and $y=x^{2}$ & Polar-coordinate squeeze \\
\end{recog}
